%% ma: L11-B10-0001
%% lop: 11
%% chuong: 4
%% bai: 10
%% dang: 11.10.1
%% loai: TN4
%% muc: NB
%% dap_an: "D"
%% nguon: "de-cuong-giua-ky-1-toan-11-nam-2026-2027-truong-thpt-hoang-van-thu-ha-noi-2-8.pdf (M08, Đề cương giữa kì 1 Toán 11 THPT Hoàng Văn Thụ, Hà Nội), Phần 1 Câu 22"
%% kien_thuc: "Bài 10 (đường thẳng và mặt phẳng trong không gian)"
%% trang_thai: cho-duyet
%% ly_do: "Dạng mới đề xuất 11.10.1"
\begin{ex}
Các yếu tố nào sau đây xác định một mặt phẳng duy nhất?
\choice{Bốn điểm phân biệt}{Ba điểm phân biệt}{Một điểm và một đường thẳng}{\True Hai đường thẳng cắt nhau}
\loigiai{Hai đường thẳng cắt nhau xác định duy nhất một mặt phẳng. Ba điểm phân biệt có thể thẳng hàng; điểm có thể nằm trên đường thẳng; bốn điểm có thể không đồng phẳng.}
\end{ex}
